\section{信頼性と故障}

\subsection{平成27年度 問7}
\label{q:h27-reg-apt-07}

\noindent\textbf{原典：} sources/H27/aptitude.pdf、PDF物理7頁。

\textbf{公式正答：} 1（single）。

\par\textit{原典本文の校正と独自解説は未作成。分類のみの登録項目。}

\section{工程管理・PERT}

\subsection{令和元年度 問10}
\label{q:r01-reg-apt-10}

\noindent\textbf{原典：} sources/R01/aptitude.pdf、PDF物理11頁。

\textbf{公式正答：} 3（single）。

\par\textit{原典本文の校正と独自解説は未作成。分類のみの登録項目。}

\section{最適化・意思決定}

\subsection{平成25年度 問14}
\label{q:h25-reg-apt-14}

\noindent\textbf{原典：} sources/H25/aptitude.pdf、PDF物理15頁。

\textbf{公式正答：} 1（single）。

\par\textit{原典本文の校正と独自解説は未作成。分類のみの登録項目。}

\section{暗号・認証・攻撃対策}

\subsection{平成28年度 問13}
\label{q:h28-reg-apt-13}

\noindent\textbf{原典：} sources/H28/aptitude.pdf、PDF物理13頁。

\textbf{公式正答：} 4（single）。

\par\textit{原典本文の校正と独自解説は未作成。分類のみの登録項目。}

\subsection{平成29年度 問6}
\label{q:h29-reg-apt-06}

\noindent\textbf{原典：} sources/H29/aptitude.pdf、PDF物理6頁。

\textbf{公式正答：} 1（single）。

\par\textit{原典本文の校正と独自解説は未作成。分類のみの登録項目。}

\section{生物・遺伝}

\subsection{令和2年度 問14}
\label{q:r02-reg-apt-14}

\noindent\textbf{原典：} sources/R02/aptitude.pdf、PDF物理14頁。

\textbf{公式正答：} 1（single）。

\par\textit{原典本文の校正と独自解説は未作成。分類のみの登録項目。}

\section{環境法・循環・持続可能性}

\subsection{令和2年度 問10}
\label{q:r02-reg-apt-10}

\noindent\textbf{原典：} sources/R02/aptitude.pdf、PDF物理10頁。

\textbf{公式正答：} 1（single）。

\par\textit{原典本文の校正と独自解説は未作成。分類のみの登録項目。}

\subsection{令和6年度 問11}
\label{q:r06-reg-apt-11}

\noindent\textbf{原典：} sources/R06/aptitude.pdf、PDF物理11頁。

\textbf{公式正答：} 5（single）。

\par\textit{原典本文の校正と独自解説は未作成。分類のみの登録項目。}

\section{公益・説明責任}

\subsection{平成25年度 問2}
\label{q:h25-reg-apt-02}

\noindent\textbf{原典：} sources/H25/aptitude.pdf、PDF物理3頁。

\textbf{公式正答：} 3（single）。

\par\textit{原典本文の校正と独自解説は未作成。分類のみの登録項目。}

\subsection{令和3年度 問3}
\label{q:r03-reg-apt-03}

\noindent\textbf{原典：} sources/R03/aptitude.pdf、PDF物理3頁。

\textbf{公式正答：} 1（single）。

\par\textbf{検討未確定：} アの受入努力は無条件の同意ではない。雇用者へCOI開示する技術者規範の当時版との最終照合は査読対象に残す。

\paragraph{問題の要旨} 説明責任・利益相反・研究公正・守秘義務の両立を判定する。

\Needspace{4\baselineskip}\paragraph{確認した条件・記述（独自要約）}\mbox{}\par

\noindent\textbf{原典確認：} 問題幹と各記述の主体・否定・全称・表の5候補を保存原典画像で確認。\par

\noindent\textbf{制度時点：} R03年度の原典表現を保持。現行資料の条番号や仕様を過去へ補わない。\par

\noindent\textbf{各記述の判断：} \{'ア': '○：受け手も理解のため努力する双方向の対話である。', 'イ': '○：利益相反の可能性を雇用者等へ開示する。', 'ウ': '○：研究実施の説明責任は共通価値である。', 'エ': '○：守秘と説明の双方を満たす方法を考える。'\}\par

\noindent\textbf{前提：} 説明責任は理由・根拠を相手が理解できる形で示すことである。研究領域を越える説明責任の価値と、利益相反情報の開示を確認する。守秘情報を無制限に公開する義務ではない。\par

\paragraph{問われる核心} 説明責任・利益相反・研究公正・守秘義務の両立を判定する。

\paragraph{必要知識} 説明責任は理由・根拠を相手が理解できる形で示すことである。研究領域を越える説明責任の価値と、利益相反情報の開示を確認する。守秘情報を無制限に公開する義務ではない。

\paragraph{考え方・正答を選ぶ理由（査読待ち初稿）} 四記述とも○、1である。アは受け手へ理由を吟味する姿勢を求める趣旨で、説明者の責務を受け手へ転嫁したり結論への無条件の賛成を要求したりする意味ではない。機密部分を抽象化し、権限のある相手や監督者へ説明することで守秘との両立を図る。

\Needspace{4\baselineskip}\paragraph{選択肢の検討（査読待ち初稿）}\mbox{}\par

\noindent\textbf{1：適切} ○・○・○・○。各記述の判断と全て一致する。\par

\noindent\textbf{2：不適切} ×・○・○・○。記述アの真偽が逆である。\par

\noindent\textbf{3：不適切} ○・×・○・○。記述イの真偽が逆である。\par

\noindent\textbf{4：不適切} ○・○・×・○。記述ウの真偽が逆である。\par

\noindent\textbf{5：不適切} ○・○・○・×。記述エの真偽が逆である。\par

\paragraph{再利用する考え方} 相手・権限・範囲を指定して説明し、情報公開と説明責任を同一視しない。

\paragraph{共通理論} \hyperref[theory:research:coi]{第\ref*{theory:research:coi}節}, \hyperref[theory:information:accountability]{第\ref*{theory:information:accountability}節}

\par 参考：\cite{singapore2010,scj_code2013}。

\par 難易度 3（校準初稿）、優先度未評価。

\section{安全・リスク評価}

\subsection{令和4年度 問4}
\label{q:r04-reg-apt-04}

\noindent\textbf{原典：} sources/R04/aptitude.pdf、PDF物理3頁。

\textbf{公式正答：} 1（single）。

\par\textit{原典本文の校正と独自解説は未作成。分類のみの登録項目。}

\subsection{令和7年度 問9}
\label{q:r07-reg-apt-09}

\noindent\textbf{原典：} sources/R07/aptitude.pdf、PDF物理9頁。

\textbf{公式正答：} 4（single）。

\par\textit{原典本文の校正と独自解説は未作成。分類のみの登録項目。}

\section{研究・情報倫理}

\subsection{平成23年度 問2}
\label{q:h23-reg-apt-02}

\noindent\textbf{原典：} sources/H23/aptitude.pdf、PDF物理3頁。

\textbf{公式正答：} 3（single）。

\paragraph{問題の要旨} 2006年の科学者行動規範から、社会への説明を理由としたデータ改変の不適切さを判断する。

\Needspace{4\baselineskip}\paragraph{確認した条件・記述（独自要約）}\mbox{}\par

\noindent\textbf{原典確認：} 問題幹と各記述の主体・否定・全称・表の5候補を保存原典画像で確認。\par

\noindent\textbf{制度時点：} H23年度の原典表現を保持。現行資料の条番号や仕様を過去へ補わない。\par

\noindent\textbf{各記述の判断：} \{'規範': '2006年声明の2・4・5・9・11項に照合。'\}\par

\noindent\textbf{前提：} 科学者の自律性は社会の信頼に支えられ、正直な判断と厳格な記録を必要とする。分かりやすい表現と結果自体の改変は違う。\par

\paragraph{問われる核心} 2006年の科学者行動規範から、社会への説明を理由としたデータ改変の不適切さを判断する。

\paragraph{必要知識} 科学者の自律性は社会の信頼に支えられ、正直な判断と厳格な記録を必要とする。分かりやすい表現と結果自体の改変は違う。

\paragraph{考え方・正答を選ぶ理由（査読待ち初稿）} 3の前半の厳格な記録は適切だが、説明のためにデータを修正してよいという後半が改ざん禁止に反する。他の四候補は正直・中立・相互批判・利益相反への対応であり、2006年規範に対応する。

\Needspace{4\baselineskip}\paragraph{選択肢の検討（査読待ち初稿）}\mbox{}\par

\noindent\textbf{1：適切} 正直な行動と科学者相互の評価。規範2の正直・誠実と相互評価に対応する。\par

\noindent\textbf{2：適切} 影響を中立・客観的に公開する。規範4の説明と公開に対応する。\par

\noindent\textbf{3：不適切} 理解を得るため多少のデータ修正を認める。規範5の改ざん禁止と矛盾する。説明用の図の整理や誤記訂正を、結論を変える改ざんへ拡張しない。\par

\noindent\textbf{4：適切} 批判に耳を傾け他者の成果を尊重する。規範9の他者との関係に対応する。\par

\noindent\textbf{5：適切} 利益の衝突に注意して公共性へ配慮する。規範11の利益相反への対応に合う。\par

\paragraph{再利用する考え方} 説明を受けやすくするための表現改善と、根拠を不真正にする操作を分ける。

\paragraph{共通理論} \hyperref[theory:research:norms]{第\ref*{theory:research:norms}節}, \hyperref[theory:research:ffp]{第\ref*{theory:research:ffp}節}

\par 参考：\cite{scj_code2006}。

\par 難易度 3（校準初稿）、優先度未評価。

\subsection{平成23年度 問3}
\label{q:h23-reg-apt-03}

\noindent\textbf{原典：} sources/H23/aptitude.pdf、PDF物理3頁。

\textbf{公式正答：} 5（single）。

\par\textbf{検討未確定：} 原典注記の刑法233条は2011年の懲役表現を保持する。具体的な事件の判決と現行刑罰条文全体は未照合であり、逮捕から有罪を補わない。

\paragraph{問題の要旨} 通信技術を用いた入試不正を、公正な比較条件と倫理・法の関係から評価する。

\Needspace{4\baselineskip}\paragraph{確認した条件・記述（独自要約）}\mbox{}\par

\noindent\textbf{原典確認：} 問題幹と各記述の主体・否定・全称・表の5候補を保存原典画像で確認。\par

\noindent\textbf{制度時点：} H23年度の原典表現を保持。現行資料の条番号や仕様を過去へ補わない。\par

\noindent\textbf{各記述の判断：} \{'ア': '○：共通条件の公正を維持する必要がある。', 'イ': '×：業務や信用に対する犯罪もあり得る。', 'ウ': '×：本人がリスクを負うだけでは他者への不公正を解消しない。', 'エ': '○：新技術への法整備と倫理観の育成が必要である。'\}\par

\noindent\textbf{前提：} 試験の成績は共通条件で比較する指標である。条件を破った得点は比較可能性を壊す。法が保護する利益は身体や物に限られない。\par

\paragraph{問われる核心} 通信技術を用いた入試不正を、公正な比較条件と倫理・法の関係から評価する。

\paragraph{必要知識} 試験の成績は共通条件で比較する指標である。条件を破った得点は比較可能性を壊す。法が保護する利益は身体や物に限られない。

\paragraph{考え方・正答を選ぶ理由（査読待ち初稿）} ア・エが○、イ・ウが×である。身体侵害や窃盗がないというだけでは犯罪成立の可能性は除外されず、自分が発覚リスクを引き受けても他の受験者の公正な選考は損なわれる。○××○、5である。逮捕を有罪判決と同一視しない。

\Needspace{4\baselineskip}\paragraph{選択肢の検討（査読待ち初稿）}\mbox{}\par

\noindent\textbf{1：不適切} ×・○・○・×。記述ア・イ・ウ・エの真偽が逆である。\par

\noindent\textbf{2：不適切} ○・○・×・○。記述イの真偽が逆である。\par

\noindent\textbf{3：不適切} ○・○・○・○。記述イ・ウの真偽が逆である。\par

\noindent\textbf{4：不適切} ○・×・○・×。記述ウ・エの真偽が逆である。\par

\noindent\textbf{5：適切} ○・×・×・○。各記述の判断と全て一致する。\par

\paragraph{再利用する考え方} 新技術が行為を可能にしたことから、その行為の許容性を推論しない。

\paragraph{共通理論} \hyperref[theory:information:fairness]{第\ref*{theory:information:fairness}節}

\par 参考：\cite{scj_code2006}。

\par 難易度 3（校準初稿）、優先度未評価。

\subsection{平成24年度 問15}
\label{q:h24-reg-apt-15}

\noindent\textbf{原典：} sources/H24/aptitude.pdf、PDF物理13頁。

\textbf{公式正答：} 5（single）。

\par\textit{原典本文の校正と独自解説は未作成。分類のみの登録項目。}

\subsection{平成26年度 問12}
\label{q:h26-reg-apt-12}

\noindent\textbf{原典：} sources/H26/aptitude.pdf、PDF物理12頁。

\textbf{公式正答：} 2（single）。

\par\textit{原典本文の校正と独自解説は未作成。分類のみの登録項目。}

\subsection{平成26年度 問2}
\label{q:h26-reg-apt-02}

\noindent\textbf{原典：} sources/H26/aptitude.pdf、PDF物理2頁。

\textbf{公式正答：} 5（single）。

\paragraph{問題の要旨} 2013年改訂の科学者の行動規範から、最も不適切な記述を選ぶ。

\Needspace{4\baselineskip}\paragraph{確認した条件・記述（独自要約）}\mbox{}\par

\noindent\textbf{時点：} 2006年制定、2013年改訂の日本学術会議声明。\par

\noindent\textbf{方向：} 不適切な記述を1つ選ぶ。\par

\paragraph{問われる核心} 科学の自律性と社会への責任

\paragraph{必要知識} 科学者の範囲、研究の功績と責任、説明責任、社会との対話。

\paragraph{考え方・正答を選ぶ理由（査読待ち初稿）} 各肢の主語と責任を確認する。1〜4は科学者の範囲、品質責任、説明、研究公正という規範の要素に対応する。5は科学の自律性を社会からの孤立と読み替え、対話と助言を否定しているため選ぶ。

\Needspace{4\baselineskip}\paragraph{選択肢の検討（査読待ち初稿）}\mbox{}\par

\noindent\textbf{1：問の条件に一致しない候補} 科学者の範囲は所属に依存せず人文・社会・自然科学の研究者と知識を活用する専門職を含む。前文での科学者の定義に合う。自然科学者や大学所属者に限定しない。\par

\noindent\textbf{2：問の条件に一致しない候補} 知識・技術の質を保証し、健康・福祉・安全・地球環境へ貢献する責任を持つ。基本的責任の規定に合う。専門知識の質と社会への貢献の双方を求める。\par

\noindent\textbf{3：問の条件に一致しない候補} 研究の意義と社会・環境への影響を中立・客観的に説明・公表する。説明と公開の責務に合う。専門性を理由に説明責任を免れるものではない。\par

\noindent\textbf{4：問の条件に一致しない候補} 役割に応じた功績と責任を負い、データを保存してFFPを行わず加担しない。研究活動の公正さに関する責務に合う。功績のみを受け取る立場ではない。\par

\noindent\textbf{5：問の条件に一致する候補} 独立性のため市民対話に否定的であり、政策への科学的助言も避ける。独立した科学的判断は対話や助言を避ける理由にならない。規範は市民との交流と政策への助言を促すので不適切である。\par

\paragraph{再利用する考え方} 制度の対象・定義と倫理上の責任を区別し、科学の自律性と社会への責任を具体的行為に適用する。

\paragraph{共通理論} \hyperref[theory:research:norms]{第\ref*{theory:research:norms}節}

\par 参考：\cite{scj_code2013}。

\par 難易度 2（校準初稿）、優先度未評価。

\subsection{平成26年度 問3}
\label{q:h26-reg-apt-03}

\noindent\textbf{原典：} sources/H26/aptitude.pdf、PDF物理3頁。

\textbf{公式正答：} 1（single）。

\paragraph{問題の要旨} 捏造・改ざん・盗用という3種の研究不正の定義を判定する。

\Needspace{4\baselineskip}\paragraph{確認した条件・記述（独自要約）}\mbox{}\par

\noindent\textbf{判断順：} 記述アから順に○=適切、×=不適切。\par

\noindent\textbf{ア：} 捏造は存在しないデータや研究結果を作ること。\par

\noindent\textbf{イ：} 改ざんは資料・機器・過程などを操作し、データや結果を真正でないものに加工すること。\par

\noindent\textbf{ウ：} 盗用は他者のアイデア・分析法・データ・結果・論文・用語を、了解若しくは適切な表示なしに流用すること。\par

\paragraph{問われる核心} 研究不正FFPの定義

\paragraph{必要知識} 捏造・改ざん・盗用の対象、データ処理と不真正化の区別。

\paragraph{考え方・正答を選ぶ理由（査読待ち初稿）} アは架空のデータの作成、イは取得・処理過程を操作した不真正化、ウは他者の成果の不適切な流用という各定義に合う。3記述とも○なので第1候補である。引用を示す責任と、著作物利用の許諾の問題は別に検討し、単に許諾があれば他者の成果を自分のものと表示してよいとは考えない。

\Needspace{4\baselineskip}\paragraph{選択肢の検討（査読待ち初稿）}\mbox{}\par

\noindent\textbf{1：問の条件に一致する候補} ○○○。全記述の判定○○○に一致する。\par

\noindent\textbf{2：問の条件に一致しない候補} ×○○。記述アの判定が逆である。独立判定は○○○となる。\par

\noindent\textbf{3：問の条件に一致しない候補} ○×○。記述イの判定が逆である。独立判定は○○○となる。\par

\noindent\textbf{4：問の条件に一致しない候補} ××○。記述ア・イの判定が逆である。独立判定は○○○となる。\par

\noindent\textbf{5：問の条件に一致しない候補} ×××。記述ア・イ・ウの判定が逆である。独立判定は○○○となる。\par

\paragraph{再利用する考え方} 制度の対象・定義と倫理上の責任を区別し、研究不正FFPの定義を具体的行為に適用する。

\paragraph{共通理論} \hyperref[theory:research:ffp]{第\ref*{theory:research:ffp}節}

\par 参考：\cite{mext_misconduct2014}。

\par 難易度 2（校準初稿）、優先度未評価。

\subsection{平成27年度 問8}
\label{q:h27-reg-apt-08}

\noindent\textbf{原典：} sources/H27/aptitude.pdf、PDF物理8頁。

\textbf{公式正答：} 1（single）。

\paragraph{問題の要旨} 実験ノートの記録・保存・訂正・異動時の持出しを判定する。

\Needspace{4\baselineskip}\paragraph{確認した条件・記述（独自要約）}\mbox{}\par

\noindent\textbf{判断順：} 記述アから順に○=適切、×=不適切。\par

\noindent\textbf{ア：} 原則として綴じたノートを消せない筆記具で記載する。\par

\noindent\textbf{イ：} 研究ノートは特許紛争回避や権利関係の明確化にも役立ち、企業研究者にも有用である。\par

\noindent\textbf{ウ：} 研究機関は一定期間の研究データ保存と必要時の開示を義務付ける規程を整える。\par

\noindent\textbf{エ：} 記述の誤りがあるページは廃棄し、見やすく修正するのが望ましい。\par

\noindent\textbf{オ：} 研究者個人に帰属するので異動時には個人判断で持ち出せる。\par

\paragraph{問われる核心} 研究記録の追跡可能性と機関管理

\paragraph{必要知識} 元記録の保存、訂正の履歴、データ帰属と異動時の管理。

\paragraph{考え方・正答を選ぶ理由（査読待ち初稿）} アは紙ノートの追記・差替えを追跡できるようにする原則である。イは研究の経緯を立証する記録として有用という記述であり、ノートだけで特許権が自動的に決まるとの意味ではない。ウは機関の保存規程に対応する。エのページ廃棄は元の経過を失わせる。オの個人所有・自由持出しの一律断定も不適切であり、資金・契約・機関規程による管理が必要である。○○○××の第1候補を選ぶ。

\Needspace{4\baselineskip}\paragraph{選択肢の検討（査読待ち初稿）}\mbox{}\par

\noindent\textbf{1：問の条件に一致する候補} ○○○××。全記述の判定○○○××に一致する。\par

\noindent\textbf{2：問の条件に一致しない候補} ×○○××。記述アの判定が逆である。独立判定は○○○××となる。\par

\noindent\textbf{3：問の条件に一致しない候補} ○×○××。記述イの判定が逆である。独立判定は○○○××となる。\par

\noindent\textbf{4：問の条件に一致しない候補} ×○×○×。記述ア・ウ・エの判定が逆である。独立判定は○○○××となる。\par

\noindent\textbf{5：問の条件に一致しない候補} ○○○×○。記述オの判定が逆である。独立判定は○○○××となる。\par

\paragraph{再利用する考え方} 制度の対象・定義と倫理上の責任を区別し、研究記録の追跡可能性と機関管理を具体的行為に適用する。

\paragraph{共通理論} \hyperref[theory:research:records]{第\ref*{theory:research:records}節}

\par 参考：\cite{mext_misconduct2014,jsps_integrity2025}。

\par 難易度 2（校準初稿）、優先度未評価。

\subsection{平成28年度 問4}
\label{q:h28-reg-apt-04}

\noindent\textbf{原典：} sources/H28/aptitude.pdf、PDF物理4頁。

\textbf{公式正答：} 2（single）。

\paragraph{問題の要旨} 2014年ガイドラインの特定不正行為と研究倫理全般を区別して判定する。

\Needspace{4\baselineskip}\paragraph{確認した条件・記述（独自要約）}\mbox{}\par

\noindent\textbf{判断順：} 記述アから順に○=適切、×=不適切。\par

\noindent\textbf{ア：} 故意又は基本的注意義務を著しく怠った、発表成果中の捏造・改ざん・盗用が特定不正行為である。\par

\noindent\textbf{イ：} 確認された不正に係る競争的資金について配分機関は事案に応じ交付取消しや一部・全部の返還を求められる。\par

\noindent\textbf{ウ：} 二重投稿や不適切な著者表示は研究者倫理に反するが不正行為ではない。\par

\noindent\textbf{エ：} 不正の対応は研究者・コミュニティ・機関の規律と自浄作用に基づく。\par

\noindent\textbf{オ：} 研究倫理教育は倫理を向上させるが、不正予防の環境整備にはならない。\par

\paragraph{問われる核心} 特定不正行為と広い研究倫理違反

\paragraph{必要知識} FFPの限定、一般的な不正、競争的資金の措置、研究倫理教育。

\paragraph{考え方・正答を選ぶ理由（査読待ち初稿）} アは特定不正行為の要件である。イは対象の競争的資金と事案に応じた措置に限る記述で妥当であり、全資金の自動返還ではない。ウは特定不正行為の限定を不正一般へ広げている。二重投稿等も倫理に反する不正として認識されるので×である。エは自己規律・自律の位置付けに合い、オは教育の予防的役割を否定するため×である。○○×○×の第2候補となる。

\Needspace{4\baselineskip}\paragraph{選択肢の検討（査読待ち初稿）}\mbox{}\par

\noindent\textbf{1：問の条件に一致しない候補} ○○○○×。記述ウの判定が逆である。独立判定は○○×○×となる。\par

\noindent\textbf{2：問の条件に一致する候補} ○○×○×。全記述の判定○○×○×に一致する。\par

\noindent\textbf{3：問の条件に一致しない候補} ○×○×○。記述イ・ウ・エ・オの判定が逆である。独立判定は○○×○×となる。\par

\noindent\textbf{4：問の条件に一致しない候補} ×○×○×。記述アの判定が逆である。独立判定は○○×○×となる。\par

\noindent\textbf{5：問の条件に一致しない候補} ○×○××。記述イ・ウ・エの判定が逆である。独立判定は○○×○×となる。\par

\paragraph{再利用する考え方} 制度の対象・定義と倫理上の責任を区別し、特定不正行為と広い研究倫理違反を具体的行為に適用する。

\paragraph{共通理論} \hyperref[theory:research:ffp]{第\ref*{theory:research:ffp}節}, \hyperref[theory:research:prevention]{第\ref*{theory:research:prevention}節}

\par 参考：\cite{mext_misconduct2014}。

\par 難易度 2（校準初稿）、優先度未評価。

\subsection{平成29年度 問14}
\label{q:h29-reg-apt-14}

\noindent\textbf{原典：} sources/H29/aptitude.pdf、PDF物理13頁。

\textbf{公式正答：} 3（single）。

\paragraph{問題の要旨} 研究の自由、記録、資金、著者責任、引用に関する7記述中の不適切数を求める。

\Needspace{4\baselineskip}\paragraph{確認した条件・記述（独自要約）}\mbox{}\par

\noindent\textbf{方向：} 不適切な記述の個数。候補1〜5はそれぞれ0〜4個。\par

\noindent\textbf{ア：} 研究の自由は真理又は公益、独創性と正確さ、結果への責任を伴う。\par

\noindent\textbf{イ：} 客観的検証のため記録を保存し、FFPを禁じる。\par

\noindent\textbf{ウ：} 研究費は計画・申請に基づき公正に使い、目的外利用・不正操作をしない。\par

\noindent\textbf{エ：} 貢献者を著者とし論文の正確さなどへ責任を負うが、共著者は自分に関係した部分だけへ責任を持てばよい。\par

\noindent\textbf{オ：} 実験の大きな貢献があれば、論文の内容と正確さを説明できなくても共著者になれる。\par

\noindent\textbf{カ：} 先行研究を明示し、自らの独創的寄与を区別する。\par

\noindent\textbf{キ：} 著作物の引用部分と原著者を明示し、引用の範囲を超えるコピーには原則として許諾を得る。\par

\paragraph{問われる核心} 著者の貢献と説明責任

\paragraph{必要知識} 功績の帰属、共同研究での責任、引用と著作物利用の区別。

\paragraph{考え方・正答を選ぶ理由（査読待ち初稿）} ア・イ・ウは研究の目的、検証、資金責任、カ・キは先行成果と利用条件の明示で適切である。エは一律に自分の部分だけへ責任を限定しており不適切である。役割分担を明示する場合も、共同研究の信頼性の確認を放棄することにはならない。オは実験への貢献だけで内容を説明できない人を著者にできるとする点が不適切である。誤りは2件、第3候補となる。

\Needspace{4\baselineskip}\paragraph{選択肢の検討（査読待ち初稿）}\mbox{}\par

\noindent\textbf{1：問の条件に一致しない候補} 0個。エ・オの2つの誤りを数えていない。\par

\noindent\textbf{2：問の条件に一致しない候補} 1個。エの責任限定とオの説明不能の2つが不適切である。\par

\noindent\textbf{3：問の条件に一致する候補} 2個。エ・オの2記述を不適切と数えた結果に一致する。\par

\noindent\textbf{4：問の条件に一致しない候補} 3個。ア・イ・ウ・カ・キには、記載された倫理原則に反する内容はない。\par

\noindent\textbf{5：問の条件に一致しない候補} 4個。不適切はエ・オだけであり、4件ではない。\par

\paragraph{再利用する考え方} 制度の対象・定義と倫理上の責任を区別し、著者の貢献と説明責任を具体的行為に適用する。

\paragraph{共通理論} \hyperref[theory:research:authorship]{第\ref*{theory:research:authorship}節}

\par 参考：\cite{jsps_integrity2025}。

\par 難易度 2（校準初稿）、優先度未評価。

\subsection{令和元年度再試験 問10}
\label{q:r01-ret-apt-10}

\noindent\textbf{原典：} sources/R01\_retest/aptitude.pdf、PDF物理10頁。

\textbf{公式正答：} 2（single）。

\paragraph{問題の要旨} 誠実な研究の誤り、私費研究、言語の異なる重複発表、共著論文の翻訳利用を判定する。

\Needspace{4\baselineskip}\paragraph{確認した条件・記述（独自要約）}\mbox{}\par

\noindent\textbf{判断順：} 記述アから順に○=適切、×=不適切。\par

\noindent\textbf{ア：} 適切な方法で正当に得た研究成果が結果的に誤りでも、ガイドライン公表後は不正行為となった。\par

\noindent\textbf{イ：} 個人の資金と労力で研究・公表するならガイドラインの内容を考慮する必要がない。\par

\noindent\textbf{ウ：} 同一成果を日本語と英語で別の学会に発表すれば二重投稿にはならない。\par

\noindent\textbf{エ：} AとBの共著英語論文の一部を、Aが単独日本語書籍に翻訳利用する前にBへ相談して了解を得た。\par

\paragraph{問われる核心} 研究の誠実さと公表時の透明性

\paragraph{必要知識} 誠実な誤りの除外、資金源に依存しない倫理、重複発表と共同成果利用。

\paragraph{考え方・正答を選ぶ理由（査読待ち初稿）} アは誠実な誤りを不正とする点、イは資金源だけで倫理責任を免れる点、ウは言語だけで重複発表を正当化する点がいずれも誤りである。エは共同成果の再利用を独断で行わず共著者へ相談した行為として適切である。出版側の許諾・引用表示などが不要との記述ではなく、実利用時には別途確認する。×××○の第2候補である。

\Needspace{4\baselineskip}\paragraph{選択肢の検討（査読待ち初稿）}\mbox{}\par

\noindent\textbf{1：問の条件に一致しない候補} ×○×○。記述イの判定が逆である。独立判定は×××○となる。\par

\noindent\textbf{2：問の条件に一致する候補} ×××○。全記述の判定×××○に一致する。\par

\noindent\textbf{3：問の条件に一致しない候補} ○××○。記述アの判定が逆である。独立判定は×××○となる。\par

\noindent\textbf{4：問の条件に一致しない候補} ○○○×。記述ア・イ・ウ・エの判定が逆である。独立判定は×××○となる。\par

\noindent\textbf{5：問の条件に一致しない候補} ××○○。記述ウの判定が逆である。独立判定は×××○となる。\par

\paragraph{再利用する考え方} 制度の対象・定義と倫理上の責任を区別し、研究の誠実さと公表時の透明性を具体的行為に適用する。

\paragraph{共通理論} \hyperref[theory:research:ffp]{第\ref*{theory:research:ffp}節}, \hyperref[theory:research:authorship]{第\ref*{theory:research:authorship}節}

\par 参考：\cite{mext_misconduct2014,jsps_integrity2025}。

\par 難易度 2（校準初稿）、優先度未評価。

\subsection{令和元年度再試験 問6}
\label{q:r01-ret-apt-06}

\noindent\textbf{原典：} sources/R01\_retest/aptitude.pdf、PDF物理6頁。

\textbf{公式正答：} 4（single）。

\paragraph{問題の要旨} 科学者の範囲・誠実さ・両義的利用・研究環境の責任を判定する。

\Needspace{4\baselineskip}\paragraph{確認した条件・記述（独自要約）}\mbox{}\par

\noindent\textbf{判断順：} 記述アから順に○=適切、×=不適切。\par

\noindent\textbf{ア：} 科学者は所属にかかわらず全学術分野の研究者や科学知識の利活用を行う専門職を含む。\par

\noindent\textbf{イ：} 正直・誠実に行動し、専門能力の維持向上と知識の正確さの提示に努める。\par

\noindent\textbf{ウ：} 意図に反し悪用され得ても、社会へ貢献すると判断すれば速やかに研究・公表を積極的に行う。\par

\noindent\textbf{エ：} 公正な研究環境の維持と不正抑止の教育啓発に継続して取り組む責任がある。\par

\paragraph{問われる核心} 科学研究の利用の両義性

\paragraph{必要知識} 研究の便益と悪用リスク、社会に許容される研究・公表方法。

\paragraph{考え方・正答を選ぶ理由（査読待ち初稿）} アは前文の定義、イは誠実な姿勢、エは環境整備と教育の責任に対応する。ウは善用の可能性だけで悪用への配慮を省いている。規範は研究と公表の際に社会へ許容される手段・方法を選ぶことを求めるので×である。○○×○の第4候補である。

\Needspace{4\baselineskip}\paragraph{選択肢の検討（査読待ち初稿）}\mbox{}\par

\noindent\textbf{1：問の条件に一致しない候補} ○○○○。記述ウの判定が逆である。独立判定は○○×○となる。\par

\noindent\textbf{2：問の条件に一致しない候補} ×○○○。記述ア・ウの判定が逆である。独立判定は○○×○となる。\par

\noindent\textbf{3：問の条件に一致しない候補} ○×○○。記述イ・ウの判定が逆である。独立判定は○○×○となる。\par

\noindent\textbf{4：問の条件に一致する候補} ○○×○。全記述の判定○○×○に一致する。\par

\noindent\textbf{5：問の条件に一致しない候補} ○○○×。記述ウ・エの判定が逆である。独立判定は○○×○となる。\par

\paragraph{再利用する考え方} 制度の対象・定義と倫理上の責任を区別し、科学研究の利用の両義性を具体的行為に適用する。

\paragraph{共通理論} \hyperref[theory:research:norms]{第\ref*{theory:research:norms}節}, \hyperref[theory:research:prevention]{第\ref*{theory:research:prevention}節}

\par 参考：\cite{scj_code2013}。

\par 難易度 2（校準初稿）、優先度未評価。

\subsection{令和2年度 問3}
\label{q:r02-reg-apt-03}

\noindent\textbf{原典：} sources/R02/aptitude.pdf、PDF物理3頁。

\textbf{公式正答：} 3（single）。

\paragraph{問題の要旨} 兼業変更、競争相手の査読、配偶者名義の株式、公的設備の企業利用の利益相反管理を判定する。

\Needspace{4\baselineskip}\paragraph{確認した条件・記述（独自要約）}\mbox{}\par

\noindent\textbf{判断順：} 記述アから順に○=適切、×=不適切。\par

\noindent\textbf{ア：} 水曜日に許可された有償兼業について土曜の技術会議参加が必要となり、大学COI委員会へ相談した。\par

\noindent\textbf{イ：} 自己の研究と競争関係にある論文の公正な査読に不安があるため辞退した。\par

\noindent\textbf{ウ：} 公的資金で評価した企業の未公開株を少額で配偶者名義にて買い、COI委員会への相談を省いた。\par

\noindent\textbf{エ：} 大学の兼業許可を得て設立した企業が公的資金の研究設備を無償利用する必要が生じ、事前に大学COI委員会へ相談した。\par

\paragraph{問われる核心} 利益相反の開示と第三者管理

\paragraph{必要知識} 経済的利益と責務の衝突、家族の利害、査読の競争関係、機関資源。

\paragraph{考え方・正答を選ぶ理由（査読待ち初稿）} アは許可条件の変更を自己判断せず相談する対応、イは査読の公正さを確保する辞退で適切である。ウは家族名義や少額を理由に利害関係を隠しており不適切である。エは所属機関の資源と私企業の利益を分ける事前相談として適切である。相談すれば無償利用が必ず認められるとの意味ではない。○○×○の第3候補を選ぶ。

\Needspace{4\baselineskip}\paragraph{選択肢の検討（査読待ち初稿）}\mbox{}\par

\noindent\textbf{1：問の条件に一致しない候補} ○○○○。記述ウの判定が逆である。独立判定は○○×○となる。\par

\noindent\textbf{2：問の条件に一致しない候補} ○○○×。記述ウ・エの判定が逆である。独立判定は○○×○となる。\par

\noindent\textbf{3：問の条件に一致する候補} ○○×○。全記述の判定○○×○に一致する。\par

\noindent\textbf{4：問の条件に一致しない候補} ○×○○。記述イ・ウの判定が逆である。独立判定は○○×○となる。\par

\noindent\textbf{5：問の条件に一致しない候補} ×○○○。記述ア・ウの判定が逆である。独立判定は○○×○となる。\par

\paragraph{再利用する考え方} 制度の対象・定義と倫理上の責任を区別し、利益相反の開示と第三者管理を具体的行為に適用する。

\paragraph{共通理論} \hyperref[theory:research:coi]{第\ref*{theory:research:coi}節}

\par 参考：\cite{jsps_integrity2025}。

\par 難易度 2（校準初稿）、優先度未評価。

\subsection{令和2年度 問8}
\label{q:r02-reg-apt-08}

\noindent\textbf{原典：} sources/R02/aptitude.pdf、PDF物理8頁。

\textbf{公式正答：} 1（single）。

\par\textit{原典本文の校正と独自解説は未作成。分類のみの登録項目。}

\subsection{令和4年度 問15}
\label{q:r04-reg-apt-15}

\noindent\textbf{原典：} sources/R04/aptitude.pdf、PDF物理14頁。

\textbf{公式正答：} 5（single）。

\par\textit{原典本文の校正と独自解説は未作成。分類のみの登録項目。}

\subsection{令和4年度 問2}
\label{q:r04-reg-apt-02}

\noindent\textbf{原典：} sources/R04/aptitude.pdf、PDF物理2頁。

\textbf{公式正答：} 4（single）。

\par\textit{原典本文の校正と独自解説は未作成。分類のみの登録項目。}

\subsection{令和5年度 問7}
\label{q:r05-reg-apt-07}

\noindent\textbf{原典：} sources/R05/aptitude.pdf、PDF物理7頁。

\textbf{公式正答：} 2（single）。

\paragraph{問題の要旨} 科学者の研究公正、対話、科学的助言と政策判断の関係を判定する。

\Needspace{4\baselineskip}\paragraph{確認した条件・記述（独自要約）}\mbox{}\par

\noindent\textbf{判断順：} 記述アから順に○=適切、×=不適切。\par

\noindent\textbf{ア：} 公表は役割に応じた功績と責任を伴い、データを保存してFFPを行わず加担しない。\par

\noindent\textbf{イ：} 市民対話と政策助言に努め、科学者に意見の相違がある場合は多数決で統一見解を決めて助言する。\par

\noindent\textbf{ウ：} 客観的根拠に基づき権威を濫用せず、科学的不確実性と見解の多様性も説明する。\par

\noindent\textbf{エ：} 科学的知見は政策の唯一の根拠ではなく、助言と異なる決定には必要に応じ社会への説明を求める。\par

\paragraph{問われる核心} 科学的助言の不確実性と政策責任

\paragraph{必要知識} 合意と意見相違の説明、科学的根拠と政策決定の他の判断材料。

\paragraph{考え方・正答を選ぶ理由（査読待ち初稿）} ア・ウ・エはそれぞれ研究公正、助言の質、科学と政策の役割に合う。イの前半は適切だが、多数決で相違を消すという後半が誤りである。規範は意見が異なる場合に相違を説明することを求めている。○×○○の第2候補である。

\Needspace{4\baselineskip}\paragraph{選択肢の検討（査読待ち初稿）}\mbox{}\par

\noindent\textbf{1：問の条件に一致しない候補} ×○○○。記述ア・イの判定が逆である。独立判定は○×○○となる。\par

\noindent\textbf{2：問の条件に一致する候補} ○×○○。全記述の判定○×○○に一致する。\par

\noindent\textbf{3：問の条件に一致しない候補} ○○×○。記述イ・ウの判定が逆である。独立判定は○×○○となる。\par

\noindent\textbf{4：問の条件に一致しない候補} ○○○×。記述イ・エの判定が逆である。独立判定は○×○○となる。\par

\noindent\textbf{5：問の条件に一致しない候補} ○○○○。記述イの判定が逆である。独立判定は○×○○となる。\par

\paragraph{再利用する考え方} 制度の対象・定義と倫理上の責任を区別し、科学的助言の不確実性と政策責任を具体的行為に適用する。

\paragraph{共通理論} \hyperref[theory:research:norms]{第\ref*{theory:research:norms}節}

\par 参考：\cite{scj_code2013}。

\par 難易度 2（校準初稿）、優先度未評価。

\subsection{令和6年度 問8}
\label{q:r06-reg-apt-08}

\noindent\textbf{原典：} sources/R06/aptitude.pdf、PDF物理8頁。

\textbf{公式正答：} 3（single）。

\paragraph{問題の要旨} 贈与・幽霊著者、資金獲得のみの著者資格、細切れ出版を判定する。

\Needspace{4\baselineskip}\paragraph{確認した条件・記述（独自要約）}\mbox{}\par

\noindent\textbf{判断順：} 記述アから順に○=適切、×=不適切。\par

\noindent\textbf{ア：} 資格を満たさない者を好意や強制により著者へ入れる行為はギフト・オーサーシップという倫理違反である。\par

\noindent\textbf{イ：} 研究費を獲得した者は研究遂行へ寄与したので論文著者の資格を満たす。\par

\noindent\textbf{ウ：} 1つの研究を複数の小研究に分け細切れに公表することをサラミ又はボローニャ出版という。\par

\noindent\textbf{エ：} 実質的貢献と著者要件を満たす者を著者として表示しない行為はゴースト・オーサーシップという倫理違反である。\par

\paragraph{問われる核心} 著者資格と研究成果の適切な単位

\paragraph{必要知識} 貢献と責任、謝辞、ギフト・ゴースト著者、細切れ出版。

\paragraph{考え方・正答を選ぶ理由（査読待ち初稿）} アは不相応な著者の追加、エは資格ある著者の省略という定義に合う。イは資金獲得のみで著者資格を充足するとする誤りである。ウは細切れ公表の用語を説明しており○である。独立した科学的問いを持つ複数論文すべてが倫理違反という意味に一般化しない。○×○○の第3候補を選ぶ。

\Needspace{4\baselineskip}\paragraph{選択肢の検討（査読待ち初稿）}\mbox{}\par

\noindent\textbf{1：問の条件に一致しない候補} ○○○○。記述イの判定が逆である。独立判定は○×○○となる。\par

\noindent\textbf{2：問の条件に一致しない候補} ×○○○。記述ア・イの判定が逆である。独立判定は○×○○となる。\par

\noindent\textbf{3：問の条件に一致する候補} ○×○○。全記述の判定○×○○に一致する。\par

\noindent\textbf{4：問の条件に一致しない候補} ○○×○。記述イ・ウの判定が逆である。独立判定は○×○○となる。\par

\noindent\textbf{5：問の条件に一致しない候補} ○○○×。記述イ・エの判定が逆である。独立判定は○×○○となる。\par

\paragraph{再利用する考え方} 制度の対象・定義と倫理上の責任を区別し、著者資格と研究成果の適切な単位を具体的行為に適用する。

\paragraph{共通理論} \hyperref[theory:research:authorship]{第\ref*{theory:research:authorship}節}

\par 参考：\cite{jsps_integrity2025}。

\par 難易度 2（校準初稿）、優先度未評価。

\section{暗黙知・形式知の変換}

\subsection{平成23年度 問10}
\label{q:h23-reg-apt-10}

\noindent\textbf{原典：} sources/H23/aptitude.pdf、PDF物理11頁。

\textbf{公式正答：} 3（single）。

\par\textbf{検討未確定：} 原典のISO27000と2005年の組合せは、27000:2009公開書誌と不一致である。原典を27001へ無断修正しない。規格全文の条項は未取得。

\paragraph{問題の要旨} 暗黙知と形式知、および情報セキュリティの三要素を語句と定義から対応させる。

\Needspace{4\baselineskip}\paragraph{確認した条件・記述（独自要約）}\mbox{}\par

\noindent\textbf{原典確認：} 問題幹と各記述の主体・否定・全称・表の5候補を保存原典画像で確認。\par

\noindent\textbf{制度時点：} H23年度の原典表現を保持。現行資料の条番号や仕様を過去へ補わない。\par

\noindent\textbf{各記述の判断：} \{'対応': 'ア暗黙知・イ形式知・ウ機密性・エ完全性・オ可用性。', '規格疑義': '原典は2005年ISO27000と記すが、公開書誌の27000初版は2009年である。'\}\par

\noindent\textbf{前提：} 暗黙知は経験や技能として身に付く知、形式知は文章・図表等で表現された知である。機密性は権限、完全性は正確さ、可用性は必要時の利用を扱う。\par

\paragraph{問われる核心} 暗黙知と形式知、および情報セキュリティの三要素を語句と定義から対応させる。

\paragraph{必要知識} 暗黙知は経験や技能として身に付く知、形式知は文章・図表等で表現された知である。機密性は権限、完全性は正確さ、可用性は必要時の利用を扱う。

\paragraph{考え方・正答を選ぶ理由（査読待ち初稿）} 技能の勘所を暗黙知、マニュアルを形式知へ対応させる。認可された者だけという制限は機密性、正確で完全な状態は完全性、必要時に利用できることは可用性である。全て合う3を選ぶ。即時性は可用性の名称ではない。

\Needspace{4\baselineskip}\paragraph{選択肢の検討（査読待ち初稿）}\mbox{}\par

\noindent\textbf{1：不適切} 暗黙知・形式知・機密性・完全性・即時性。最後は可用性であり、即時性ではない。\par

\noindent\textbf{2：不適切} 経験知・客観知・特権性・完全性・即時性。知識管理の対語とCIAの名称が合わない。\par

\noindent\textbf{3：適切} 暗黙知・形式知・機密性・完全性・可用性。五つの定義に全て対応する。\par

\noindent\textbf{4：不適切} 暗黙知・客観知・機密性・保存性・可用性。形式知と完全性の位置を別語にしている。\par

\noindent\textbf{5：不適切} 経験知・形式知・特権性・保存性・可用性。暗黙知・機密性・完全性の標準名称に合わない。\par

\paragraph{再利用する考え方} 技能共有と情報保護を別の体系として整理し、用語を定義から判断する。

\paragraph{共通理論} \hyperref[theory:information:knowledge]{第\ref*{theory:information:knowledge}節}, \hyperref[theory:information:cia]{第\ref*{theory:information:cia}節}

\par 参考：\cite{jaist_knowledge,iso27000_2009}。

\paragraph{関連問題} \hyperref[q:r04-reg-apt-13]{R04-REG-APT-13}

\par 難易度 3（校準初稿）、優先度未評価。

\section{身体特徴の個人識別符号}

\subsection{令和元年度 問4}
\label{q:r01-reg-apt-04}

\noindent\textbf{原典：} sources/R01/aptitude.pdf、PDF物理5頁。

\textbf{公式正答：} 1（single）。

\par\textbf{検討未確定：} 現在の公開通則の列挙を確認。令和元年の施行令版の逐条差分までは未照合である。

\paragraph{問題の要旨} 身体特徴を変換した個人識別符号の七分類を判定する。

\Needspace{4\baselineskip}\paragraph{確認した条件・記述（独自要約）}\mbox{}\par

\noindent\textbf{原典確認：} 問題幹と各記述の主体・否定・全称・表の5候補を保存原典画像で確認。\par

\noindent\textbf{制度時点：} R01年度の原典表現を保持。現行資料の条番号や仕様を過去へ補わない。\par

\noindent\textbf{各記述の判断：} \{'ア': 'DNA', 'イ': '顔', 'ウ': '虹彩', 'エ': '声', 'オ': '歩行', 'カ': '静脈', 'キ': '指紋・掌紋'\}\par

\noindent\textbf{前提：} 政令のDNA・顔・虹彩・声・歩行・静脈・指紋掌紋に対応する。生の特徴全てではなく、電子計算機向けに変換され識別に足りる基準を満たす符号という前提である。\par

\paragraph{問われる核心} 身体特徴を変換した個人識別符号の七分類を判定する。

\paragraph{必要知識} 政令のDNA・顔・虹彩・声・歩行・静脈・指紋掌紋に対応する。生の特徴全てではなく、電子計算機向けに変換され識別に足りる基準を満たす符号という前提である。

\paragraph{考え方・正答を選ぶ理由（査読待ち初稿）} 七つの記述は政令第1条の七生体分類と全て対応する。含まれないもの0個、1である。列挙された身体そのものをそのまま符号と考えず、問題幹の変換・識別の前提を各記述に適用する。

\Needspace{4\baselineskip}\paragraph{選択肢の検討（査読待ち初稿）}\mbox{}\par

\noindent\textbf{1：適切} 0。含まれないなものは0個であり、一致する。\par

\noindent\textbf{2：不適切} 1。含まれないなものは0個であり、この個数とは一致しない。\par

\noindent\textbf{3：不適切} 2。含まれないなものは0個であり、この個数とは一致しない。\par

\noindent\textbf{4：不適切} 3。含まれないなものは0個であり、この個数とは一致しない。\par

\noindent\textbf{5：不適切} 4。含まれないなものは0個であり、この個数とは一致しない。\par

\paragraph{再利用する考え方} 生体の特徴と、認証用に処理された識別符号を区別する。

\paragraph{共通理論} \hyperref[theory:privacy:biometrics]{第\ref*{theory:privacy:biometrics}節}

\par 参考：\cite{privacy_general}。

\paragraph{関連問題} \hyperref[q:r06-reg-apt-09]{R06-REG-APT-09}

\par 難易度 3（校準初稿）、優先度未評価。

\subsection{令和6年度 問9}
\label{q:r06-reg-apt-09}

\noindent\textbf{原典：} sources/R06/aptitude.pdf、PDF物理9頁。

\textbf{公式正答：} 5（single）。

\par\textbf{検討未確定：} 出題当時の施行令版全文との差分は未照合。5の原典生理表現を肯定する説明ではなく、法令列挙の設問として扱う。

\paragraph{問題の要旨} 施行令に列挙される生体特徴から、骨の振動特性という非列挙項目を区別する。

\Needspace{4\baselineskip}\paragraph{確認した条件・記述（独自要約）}\mbox{}\par

\noindent\textbf{原典確認：} 問題幹と各記述の主体・否定・全称・表の5候補を保存原典画像で確認。\par

\noindent\textbf{制度時点：} R06年度の原典表現を保持。現行資料の条番号や仕様を過去へ補わない。\par

\noindent\textbf{各記述の判断：} \{'列挙': '1指紋掌紋・2顔・3声・4歩行は該当、5の振動特性は列挙にない。'\}\par

\noindent\textbf{前提：} 個人識別符号の生体分類は政令の列挙を確認する。生体認証に使えそうな特徴という一般論と、政令に定める分類は同じではない。\par

\paragraph{問われる核心} 施行令に列挙される生体特徴から、骨の振動特性という非列挙項目を区別する。

\paragraph{必要知識} 個人識別符号の生体分類は政令の列挙を確認する。生体認証に使えそうな特徴という一般論と、政令に定める分類は同じではない。

\paragraph{考え方・正答を選ぶ理由（査読待ち初稿）} 1から4は七生体分類の一部である。5の骨を振動させたときの特性はこの列挙にないので5である。原典の内耳を介さずという医学的表現を正しい生理記述へ校正し直して答えを作らない。

\Needspace{4\baselineskip}\paragraph{選択肢の検討（査読待ち初稿）}\mbox{}\par

\noindent\textbf{1：記載あり} 指紋又は掌紋。施行令の指紋・掌紋分類に対応する。\par

\noindent\textbf{2：記載あり} 顔の骨格・皮膚・部位配置の容貌。顔の特徴分類に対応する。\par

\noindent\textbf{3：記載あり} 発声時の声帯・声門・声道の変化。声の特徴分類に対応する。\par

\noindent\textbf{4：記載あり} 歩行の姿勢・腕の動作・歩幅等。歩行の特徴分類に対応する。\par

\noindent\textbf{5：記載なし} 骨振動による蝸牛への振動特性。政令の七生体分類にこの項目は列挙されていない。\par

\paragraph{再利用する考え方} 技術的な識別能力と法令上の列挙を区別する。

\paragraph{共通理論} \hyperref[theory:privacy:biometrics]{第\ref*{theory:privacy:biometrics}節}

\par 参考：\cite{privacy_general}。

\paragraph{関連問題} \hyperref[q:r01-reg-apt-04]{R01-REG-APT-04}

\par 難易度 3（校準初稿）、優先度未評価。

\section{本人同意・第三者提供・利用目的}

\subsection{平成26年度 問4}
\label{q:h26-reg-apt-04}

\noindent\textbf{原典：} sources/H26/aptitude.pdf、PDF物理4頁。

\textbf{公式正答：} 4（single）。

\par\textbf{検討未確定：} 学校名簿の配布は同意原則の問題であり、法定例外を全て否定するものではない。2013年指針の当時版全文は未照合。

\paragraph{問題の要旨} 公開情報の保護と、災害時の名簿共有・学校名簿の作成と配布の違いを判定する。

\Needspace{4\baselineskip}\paragraph{確認した条件・記述（独自要約）}\mbox{}\par

\noindent\textbf{原典確認：} 問題幹と各記述の主体・否定・全称・表の5候補を保存原典画像で確認。\par

\noindent\textbf{制度時点：} H26年度の原典表現を保持。現行資料の条番号や仕様を過去へ補わない。\par

\noindent\textbf{各記述の判断：} \{'ア': '○：識別できる評価・映像・音声も対象である。', 'イ': '×：公知だから除外という規定ではない。', 'ウ': '○：法定災害条件と必要範囲内の名簿共有である。', 'エ': '○：通常の学校名簿配布で同意を必要とする。'\}\par

\noindent\textbf{前提：} 適正取得・目的通知、利用、第三者提供を分ける。災害発生又は発生のおそれで生命身体保護へ特に必要な場合は、避難支援に必要な限度で同意なし提供が認められる。\par

\paragraph{問われる核心} 公開情報の保護と、災害時の名簿共有・学校名簿の作成と配布の違いを判定する。

\paragraph{必要知識} 適正取得・目的通知、利用、第三者提供を分ける。災害発生又は発生のおそれで生命身体保護へ特に必要な場合は、避難支援に必要な限度で同意なし提供が認められる。

\paragraph{考え方・正答を選ぶ理由（査読待ち初稿）} 公開済みという理由で保護を除くイが×、他は通常の条件下で○となり○×○○、4である。災害時の例外を平時の名簿の自由な提供へ拡張しない。名簿を内部で作る行為と家庭等へ配布する行為では規律が異なる。

\Needspace{4\baselineskip}\paragraph{選択肢の検討（査読待ち初稿）}\mbox{}\par

\noindent\textbf{1：不適切} ○・×・×・○。記述ウの真偽が逆である。\par

\noindent\textbf{2：不適切} ×・○・○・○。記述ア・イの真偽が逆である。\par

\noindent\textbf{3：不適切} ×・○・×・○。記述ア・イ・ウの真偽が逆である。\par

\noindent\textbf{4：適切} ○・×・○・○。各記述の判断と全て一致する。\par

\noindent\textbf{5：不適切} ○・○・○・×。記述イ・エの真偽が逆である。\par

\paragraph{再利用する考え方} 取得・利用・提供を時点と目的に沿って分け、例外の必要範囲を確認する。

\paragraph{共通理論} \hyperref[theory:privacy:scope]{第\ref*{theory:privacy:scope}節}, \hyperref[theory:privacy:consent]{第\ref*{theory:privacy:consent}節}, \hyperref[theory:privacy:evacuation]{第\ref*{theory:privacy:evacuation}節}

\par 参考：\cite{privacy_general,privacy_evacuation2013,privacy_evacuation_scope}。

\paragraph{関連問題} \hyperref[q:h24-reg-apt-04]{H24-REG-APT-04}

\par 難易度 3（校準初稿）、優先度未評価。

\subsection{令和3年度 問14}
\label{q:r03-reg-apt-14}

\noindent\textbf{原典：} sources/R03/aptitude.pdf、PDF物理14頁。

\textbf{公式正答：} 5（single）。

\par\textbf{独立解答と公式正答に相違あり：} 独立解答 2。初回独立2と公式5は一致しない。黙示同意自体を一律禁止する説明は不正確である。加盟店への連絡を承諾した具体的事実、契約条件、当時の直接Q\&Aは未確認であり、公式へ合わせて条件を補わない。

\paragraph{問題の要旨} 謝罪目的の連絡先提供・カード照会・目的外の求人勧誘・CC送信を本人同意と漏えいから評価する。

\Needspace{4\baselineskip}\paragraph{確認した条件・記述（独自要約）}\mbox{}\par

\noindent\textbf{原典確認：} 問題幹と各記述の主体・否定・全称・表の5候補を保存原典画像で確認。\par

\noindent\textbf{制度時点：} R03年度の原典表現を保持。現行資料の条番号や仕様を過去へ補わない。\par

\noindent\textbf{各記述の判断：} \{'ア': '×：被害側の同意なしに加害側の謝罪目的で連絡先を渡さない。', 'イ': '公式では×：調査依頼と加盟店への第三者提供同意を同一視しない。', 'ウ': '×：ポイント目的から求人勧誘への変更には目的範囲と同意の検討が必要である。', 'エ': '×：CCは他の宛先へアドレスを開示する。'\}\par

\noindent\textbf{前提：} 本人同意は示された取扱方法への承諾の意思表示である。明示に限られないが黙示同意は具体的事実で判断する。本人の調査依頼だけで加盟店への連絡先提供を当然に認めない。\par

\paragraph{問われる核心} 謝罪目的の連絡先提供・カード照会・目的外の求人勧誘・CC送信を本人同意と漏えいから評価する。

\paragraph{必要知識} 本人同意は示された取扱方法への承諾の意思表示である。明示に限られないが黙示同意は具体的事実で判断する。本人の調査依頼だけで加盟店への連絡先提供を当然に認めない。

\paragraph{考え方・正答を選ぶ理由（査読待ち初稿）} 公式は四つとも×、5である。アは謝罪する側の事情で被害側の同意を代替している。ウは取得時の利用目的を求人へ広げ、エは顧客間でアドレスを漏らす。イの初回は調査依頼から黙示の承諾があると解して2としたが、PPCの定義では提供方法への承諾を認識できる具体的条件が必要であり、問題にその条件は示されない。公式5を保持し、初回2との相違と同意条件の疑義を査読へ残す。

\Needspace{4\baselineskip}\paragraph{選択肢の検討（査読待ち初稿）}\mbox{}\par

\noindent\textbf{1：不適切} ○・×・×・×。記述アの真偽が逆である。\par

\noindent\textbf{2：不適切} ×・○・×・×。記述イの真偽が逆である。\par

\noindent\textbf{3：不適切} ×・×・○・×。記述ウの真偽が逆である。\par

\noindent\textbf{4：不適切} ×・×・×・○。記述エの真偽が逆である。\par

\noindent\textbf{5：適切} ×・×・×・×。各記述の判断と全て一致する。\par

\paragraph{再利用する考え方} 便宜・善意・契約履行の必要性と、本人が承諾した提供範囲を別々に確認する。

\paragraph{共通理論} \hyperref[theory:privacy:consent]{第\ref*{theory:privacy:consent}節}

\par 参考：\cite{privacy_general,privacy_consent2017,privacy_consent_current}。

\par 難易度 3（校準初稿）、優先度未評価。

\section{個人情報の定義・公開と識別}

\subsection{平成23年度 問7}
\label{q:h23-reg-apt-07}

\noindent\textbf{原典：} sources/H23/aptitude.pdf、PDF物理8頁。

\textbf{公式正答：} 5（single）。

\par\textbf{検討未確定：} 配送物の中身を取り扱わない通常条件に依存する。当時ガイドラインの該当段落と責任主体の直接照合は未完了。現行Q7-35は中身の合意・発送者の安全管理を条件としている。

\paragraph{問題の要旨} 死者の家族・公開情報・監視画像・封入媒体の運送を、個人情報と取扱委託の条件から判定する。

\Needspace{4\baselineskip}\paragraph{確認した条件・記述（独自要約）}\mbox{}\par

\noindent\textbf{原典確認：} 問題幹と各記述の主体・否定・全称・表の5候補を保存原典画像で確認。\par

\noindent\textbf{制度時点：} H23年度の原典表現を保持。現行資料の条番号や仕様を過去へ補わない。\par

\noindent\textbf{各記述の判断：} \{'ア': '×：生存家族について識別できる情報は対象となる。', 'イ': '×：公開は個人情報性を消さない。', 'ウ': '×：個人を識別できる画像は対象となる。', 'エ': '×：中身を取り扱わない通常の配送だけで中身の個人データ取扱委託としない。'\}\par

\noindent\textbf{前提：} 死者自身の情報と生存する家族の情報を分ける。公開済みでも個人情報である。識別可能な映像は対象となる。宛名情報と配送物内部のCDデータは別の取扱いである。\par

\paragraph{問われる核心} 死者の家族・公開情報・監視画像・封入媒体の運送を、個人情報と取扱委託の条件から判定する。

\paragraph{必要知識} 死者自身の情報と生存する家族の情報を分ける。公開済みでも個人情報である。識別可能な映像は対象となる。宛名情報と配送物内部のCDデータは別の取扱いである。

\paragraph{考え方・正答を選ぶ理由（査読待ち初稿）} ア・イ・ウの一律除外は定義に反する。エは通常の配送で中身の取扱合意がないという条件で、中身のデータの委託先とはみなさない。四つとも×、5である。誤配に契約上等の責任が一切ないとはいえず、発送者の安全管理義務も残る。

\Needspace{4\baselineskip}\paragraph{選択肢の検討（査読待ち初稿）}\mbox{}\par

\noindent\textbf{1：不適切} ○・×・×・×。記述アの真偽が逆である。\par

\noindent\textbf{2：不適切} ×・○・×・×。記述イの真偽が逆である。\par

\noindent\textbf{3：不適切} ×・×・○・×。記述ウの真偽が逆である。\par

\noindent\textbf{4：不適切} ×・×・×・○。記述エの真偽が逆である。\par

\noindent\textbf{5：適切} ×・×・×・×。各記述の判断と全て一致する。\par

\paragraph{再利用する考え方} 情報の対象者と、事業者が実際に扱う情報の範囲を個別に決める。

\paragraph{共通理論} \hyperref[theory:privacy:scope]{第\ref*{theory:privacy:scope}節}, \hyperref[theory:privacy:delivery]{第\ref*{theory:privacy:delivery}節}

\par 参考：\cite{privacy_general,privacy_delivery}。

\paragraph{関連問題} \hyperref[q:h24-reg-apt-04]{H24-REG-APT-04}

\par 難易度 3（校準初稿）、優先度未評価。

\subsection{平成24年度 問4}
\label{q:h24-reg-apt-04}

\noindent\textbf{原典：} sources/H24/aptitude.pdf、PDF物理3頁。

\textbf{公式正答：} 2（single）。

\par\textbf{検討未確定：} 5000人基準の旧政令全文と重複人数の計数は追加照合が必要である。番号だけで一律に個人情報とするのでなく容易な照合の条件を明示する。

\paragraph{問題の要旨} 旧5000人基準と、評価情報・識別番号・防犯映像の個人情報性を判定する。

\Needspace{4\baselineskip}\paragraph{確認した条件・記述（独自要約）}\mbox{}\par

\noindent\textbf{原典確認：} 問題幹と各記述の主体・否定・全称・表の5候補を保存原典画像で確認。\par

\noindent\textbf{制度時点：} H24年度の原典表現を保持。現行資料の条番号や仕様を過去へ補わない。\par

\noindent\textbf{各記述の判断：} \{'ア': '○：判断・評価も個人に関する情報である。', 'イ': '×：個々のDBだけを数えて判断しない。', 'ウ': '○：本人へ容易に照合できる従業員・学籍番号等を想定する。', 'エ': '×：防犯目的は個人情報性の除外理由ではない。'\}\par

\noindent\textbf{前提：} 当時の事業者適用除外の人数は事業者全体で判断し、DBを分割しただけでは適用を免れない。番号は容易な照合による識別性を確認する。2017年に小規模事業者の除外は廃止された。\par

\paragraph{問われる核心} 旧5000人基準と、評価情報・識別番号・防犯映像の個人情報性を判定する。

\paragraph{必要知識} 当時の事業者適用除外の人数は事業者全体で判断し、DBを分割しただけでは適用を免れない。番号は容易な照合による識別性を確認する。2017年に小規模事業者の除外は廃止された。

\paragraph{考え方・正答を選ぶ理由（査読待ち初稿）} アとウは○である。ウは番号から本人へ容易に照合できる組織内の番号という通常条件に沿う。イは事業者全体の人数をDB単位へ分割して判断している。エは識別できる映像を目的だけで除外する。○×○×、2である。

\Needspace{4\baselineskip}\paragraph{選択肢の検討（査読待ち初稿）}\mbox{}\par

\noindent\textbf{1：不適切} ×・○・×・○。記述ア・イ・ウ・エの真偽が逆である。\par

\noindent\textbf{2：適切} ○・×・○・×。各記述の判断と全て一致する。\par

\noindent\textbf{3：不適切} ○・×・×・○。記述ウ・エの真偽が逆である。\par

\noindent\textbf{4：不適切} ○・○・×・×。記述イ・ウの真偽が逆である。\par

\noindent\textbf{5：不適切} ×・○・○・×。記述ア・イの真偽が逆である。\par

\paragraph{再利用する考え方} 適用対象の旧人数基準と、情報自体の識別性を混同しない。

\paragraph{共通理論} \hyperref[theory:privacy:scope]{第\ref*{theory:privacy:scope}節}, \hyperref[theory:privacy:history]{第\ref*{theory:privacy:history}節}

\par 参考：\cite{privacy_general,privacy_small2017}。

\paragraph{関連問題} \hyperref[q:h23-reg-apt-07]{H23-REG-APT-07}, \hyperref[q:h26-reg-apt-04]{H26-REG-APT-04}

\par 難易度 3（校準初稿）、優先度未評価。

\section{秘密情報の接近制御・相互確認と持出し防止}

\subsection{令和元年度 問9}
\label{q:r01-reg-apt-09}

\noindent\textbf{原典：} sources/R01/aptitude.pdf、PDF物理10頁。

\textbf{公式正答：} 2（single）。

\paragraph{問題の要旨} 秘密情報のアクセス、相互確認、表示、配置、暗号・閲覧制限、紙削減の6記述で不適切な数を求める。

\Needspace{4\baselineskip}\paragraph{確認した条件・記述（独自要約）}\mbox{}\par

\noindent\textbf{方向：} 不適切な記述の数。候補0/1/2/3/4。\par

\noindent\textbf{構成：} ア〜カの6項目。\par

\noindent\textbf{限定：} オは外部では閲覧できないように設定する対策。暗号自体の万能性や外部送信の許可ではない。\par

\paragraph{問われる核心} 秘密情報の接近制御と相互確認

\paragraph{必要知識} 秘密情報の接近制御と相互確認。ア○：情報の分類に応じたアクセス権で、必要のない人の接近を制限する。イ×：相互の確認や責任者のチェックを減らす単独作業の一律推奨は逆方向である。ウ○：秘密の表示で取扱者の認識を支える。エ○：資料・作業が死角になることを減らすが、権限のない者へ画面を見せてよい意味ではない。オ○：暗号と登録ID・端末等の閲覧権制限を適切に設定する対策である。暗号化だけで漏えいが絶対不可能になるという主張ではなく、無断送信を許すものでもない。カ○：電子的アクセス管理と組み合わせ、無権限者が紙に接する機会を減らす。不適切はイの1項目で2。

\paragraph{考え方・正答を選ぶ理由（査読待ち初稿）} ア○：情報の分類に応じたアクセス権で、必要のない人の接近を制限する。イ×：相互の確認や責任者のチェックを減らす単独作業の一律推奨は逆方向である。ウ○：秘密の表示で取扱者の認識を支える。エ○：資料・作業が死角になることを減らすが、権限のない者へ画面を見せてよい意味ではない。オ○：暗号と登録ID・端末等の閲覧権制限を適切に設定する対策である。暗号化だけで漏えいが絶対不可能になるという主張ではなく、無断送信を許すものでもない。カ○：電子的アクセス管理と組み合わせ、無権限者が紙に接する機会を減らす。不適切はイの1項目で2。

\Needspace{4\baselineskip}\paragraph{選択肢の検討（査読待ち初稿）}\mbox{}\par

\noindent\textbf{1：選択対象ではない} 0項目。単独作業を可能な限り推奨するイが不適切なので0ではない。\par

\noindent\textbf{2：選択対象} 1項目。イだけが不適切であり1項目となる。\par

\noindent\textbf{3：選択対象ではない} 2項目。イに加えて他の適切な対策も否定すると数え過ぎとなる。\par

\noindent\textbf{4：選択対象ではない} 3項目。アクセス権の設定・表示・配置・適切な閲覧制限・紙削減の全てを危険低減策として区別すれば3ではない。\par

\noindent\textbf{5：選択対象ではない} 4項目。適切な対策まで不適切に数えている。秘密保護と一人作業の推奨を混同しない。\par

\paragraph{再利用する考え方} 対象、条件、時点を先に確認する。秘密情報の接近制御と相互確認を具体的な事例へ適用し、技術的原因・法的要件・倫理的判断のどれを問うか区別する。

\paragraph{共通理論} \hyperref[theory:research:secret_management]{第\ref*{theory:research:secret_management}節}, \hyperref[theory:compliance:secret]{第\ref*{theory:compliance:secret}節}

\par 参考：\cite{meti_secret_handbook2024}。

\paragraph{関連問題} \hyperref[q:r05-reg-apt-02]{R05-REG-APT-02}

\par 難易度 2（校準初稿）、優先度未評価。

\subsection{令和5年度 問2}
\label{q:r05-reg-apt-02}

\noindent\textbf{原典：} sources/R05/aptitude.pdf、PDF物理2頁。

\textbf{公式正答：} 3（single）。

\paragraph{問題の要旨} アクセス権、秘密表示、単独作業、電子的コピー制限、座席配置という漏えい対策を判定する。

\Needspace{4\baselineskip}\paragraph{確認した条件・記述（独自要約）}\mbox{}\par

\noindent\textbf{方向：} 最も不適切。\par

\noindent\textbf{時点：} 令和5年。2024年版ハンドブックは現在の背景資料として区別する。\par

\noindent\textbf{範囲：} 5つの対策案。電子的制限の完全防御を主張していない。\par

\paragraph{問われる核心} 秘密保護におけるアクセス権と相互確認

\paragraph{必要知識} 秘密保護におけるアクセス権と相互確認。情報の接近制御、秘密の認識、持出しの困難化、相互チェックを組み合わせる。3だけはチェックを減らす単独作業を一律に推奨しているので不適切。複数人にアクセス権を無制限に与えることとは区別する。

\paragraph{考え方・正答を選ぶ理由（査読待ち初稿）} 情報の接近制御、秘密の認識、持出しの困難化、相互チェックを組み合わせる。3だけはチェックを減らす単独作業を一律に推奨しているので不適切。複数人にアクセス権を無制限に与えることとは区別する。

\Needspace{4\baselineskip}\paragraph{選択肢の検討（査読待ち初稿）}\mbox{}\par

\noindent\textbf{1：適切} 情報分類ごとにアクセス権の設定ルールと適切な範囲を定める。。必要性に基づくアクセス範囲の制御である。権限を無制限に広げる対策ではない。\par

\noindent\textbf{2：適切} 媒体へ秘密情報であることが分かる表示をする。。取扱者が秘密の分類と注意を認識する仕組みを支える。\par

\noindent\textbf{3：不適切} 複数人を避け、可能な限り単独で秘密情報を扱う。。必要に応じ複数人で相互確認できる体制を整え、単独が必要なら責任者が事前・事後に確認する。単独を最大化する方針は逆である。\par

\noindent\textbf{4：適切} 印刷・複製・外部メディアへのコピー等を制限する設定をする。。持出しの経路を減らす技術的対策となる。印刷禁止やコピーガードが撮影・再入力まで完全に防ぐとの保証ではない。\par

\noindent\textbf{5：適切} 作業・資料が死角になりにくい配置を検討する。。相互に作業へ注意が向く環境を作る。秘密画面をアクセス権のない人にさらすことを避けるという別条件も守る。\par

\paragraph{再利用する考え方} 対象、条件、時点を先に確認する。秘密保護におけるアクセス権と相互確認を具体的な事例へ適用し、技術的原因・法的要件・倫理的判断のどれを問うか区別する。

\paragraph{共通理論} \hyperref[theory:research:secret_management]{第\ref*{theory:research:secret_management}節}, \hyperref[theory:compliance:secret]{第\ref*{theory:compliance:secret}節}

\par 参考：\cite{meti_secret_handbook2024}。

\paragraph{関連問題} \hyperref[q:r01-reg-apt-09]{R01-REG-APT-09}

\par 難易度 2（校準初稿）、優先度未評価。

\section{CIA三要素と方針・管理体系}

\subsection{令和4年度 問13}
\label{q:r04-reg-apt-13}

\noindent\textbf{原典：} sources/R04/aptitude.pdf、PDF物理12頁。

\textbf{公式正答：} 3（single）。

\par\textbf{検討未確定：} 規格全文は未取得であり各措置を特定条項の逐語義務とは断定しない。

\paragraph{問題の要旨} CIA三要素と、方針・教育・監視・事故対応による継続改善を判定する。

\Needspace{4\baselineskip}\paragraph{確認した条件・記述（独自要約）}\mbox{}\par

\noindent\textbf{原典確認：} 問題幹と各記述の主体・否定・全称・表の5候補を保存原典画像で確認。\par

\noindent\textbf{制度時点：} R04年度の原典表現を保持。現行資料の条番号や仕様を過去へ補わない。\par

\noindent\textbf{各記述の判断：} \{'ア': '○：CIAを維持する。', 'イ': '×：これは可用性でなく完全性である。', 'ウ': '○：方針の有効性を継続的に確認する。', 'エ': '○：幹部も含む教育を実施する。', 'オ': '○：初動・分析・復旧・再発防止を方針へ反映する。'\}\par

\noindent\textbf{前提：} 可用性は必要時の利用、完全性は情報や処理の正確さと完全な状態である。方針を実施し監視・事故分析から改善するのが管理体系である。\par

\paragraph{問われる核心} CIA三要素と、方針・教育・監視・事故対応による継続改善を判定する。

\paragraph{必要知識} 可用性は必要時の利用、完全性は情報や処理の正確さと完全な状態である。方針を実施し監視・事故分析から改善するのが管理体系である。

\paragraph{考え方・正答を選ぶ理由（査読待ち初稿）} イだけ×であり、○×○○○、3である。正しい内容が書かれていても、対応する名称が可用性に入れ替わっている。ウ・エ・オは方針を文書で作って終わらせず、運用結果を点検して改善する枠組みに対応する。

\Needspace{4\baselineskip}\paragraph{選択肢の検討（査読待ち初稿）}\mbox{}\par

\noindent\textbf{1：不適切} ×・○・○・×・○。記述ア・イ・エの真偽が逆である。\par

\noindent\textbf{2：不適切} ×・×・○・○・○。記述アの真偽が逆である。\par

\noindent\textbf{3：適切} ○・×・○・○・○。各記述の判断と全て一致する。\par

\noindent\textbf{4：不適切} ○・○・×・○・×。記述イ・ウ・オの真偽が逆である。\par

\noindent\textbf{5：不適切} ○・○・×・○・○。記述イ・ウの真偽が逆である。\par

\paragraph{再利用する考え方} 名称の定義と、運用を継続改善する仕組みを別々に確かめる。

\paragraph{共通理論} \hyperref[theory:information:cia]{第\ref*{theory:information:cia}節}, \hyperref[theory:information:management]{第\ref*{theory:information:management}節}

\par 参考：\cite{iso27000_2009}。

\paragraph{関連問題} \hyperref[q:h23-reg-apt-10]{H23-REG-APT-10}, \hyperref[q:r07-reg-apt-05]{R07-REG-APT-05}

\par 難易度 3（校準初稿）、優先度未評価。

\subsection{令和7年度 問5}
\label{q:r07-reg-apt-05}

\noindent\textbf{原典：} sources/R07/aptitude.pdf、PDF物理5頁。

\textbf{公式正答：} 1（single）。

\par\textbf{検討未確定：} 27001の該当版の全条項は未取得。CIA名称と方針・体系の概念差を初稿として説明する。

\paragraph{問題の要旨} セキュリティ方針と管理体系の名称の入替えを、CIA三要素と比較する。

\Needspace{4\baselineskip}\paragraph{確認した条件・記述（独自要約）}\mbox{}\par

\noindent\textbf{原典確認：} 問題幹と各記述の主体・否定・全称・表の5候補を保存原典画像で確認。\par

\noindent\textbf{制度時点：} R07年度の原典表現を保持。現行資料の条番号や仕様を過去へ補わない。\par

\noindent\textbf{各記述の判断：} \{'ア': '×：ISO/IEC27001の管理体系を方針という名称にしている。', 'イ': '○：必要時の認可利用は可用性である。', 'ウ': '○：正確さと完全な状態は完全性である。', 'エ': '○：認可者だけのアクセスは機密性である。', 'オ': '×：方針の定義を管理体系という名称にしている。'\}\par

\noindent\textbf{前提：} ポリシーは組織の方針と行動指針、ISMSはその方針を計画・実施・評価・改善する管理体系である。CIAの定義と独立に名称の対応を確認する。\par

\paragraph{問われる核心} セキュリティ方針と管理体系の名称の入替えを、CIA三要素と比較する。

\paragraph{必要知識} ポリシーは組織の方針と行動指針、ISMSはその方針を計画・実施・評価・改善する管理体系である。CIAの定義と独立に名称の対応を確認する。

\paragraph{考え方・正答を選ぶ理由（査読待ち初稿）} イ・ウ・エは各三要素の定義に合う。アとオは管理体系と方針の名称を入れ替えた説明であるため×。×○○○×、1である。方針は管理体系に含まれるが、方針文書だけで管理体系の全てを代替しない。

\Needspace{4\baselineskip}\paragraph{選択肢の検討（査読待ち初稿）}\mbox{}\par

\noindent\textbf{1：適切} ×・○・○・○・×。各記述の判断と全て一致する。\par

\noindent\textbf{2：不適切} ○・×・○・○・○。記述ア・イ・オの真偽が逆である。\par

\noindent\textbf{3：不適切} ○・○・×・×・○。記述ア・ウ・エ・オの真偽が逆である。\par

\noindent\textbf{4：不適切} ×・○・○・×・○。記述エ・オの真偽が逆である。\par

\noindent\textbf{5：不適切} ○・○・○・○・×。記述アの真偽が逆である。\par

\paragraph{再利用する考え方} 同じ情報管理に関する語でも方針・実施・評価を分けて読む。

\paragraph{共通理論} \hyperref[theory:information:cia]{第\ref*{theory:information:cia}節}, \hyperref[theory:information:management]{第\ref*{theory:information:management}節}

\par 参考：\cite{iso27000_2009}。

\paragraph{関連問題} \hyperref[q:r04-reg-apt-13]{R04-REG-APT-13}

\par 難易度 3（校準初稿）、優先度未評価。

\section{社会的責任・安全の標準}

\subsection{平成29年度 問11}
\label{q:h29-reg-apt-11}

\noindent\textbf{原典：} sources/H29/aptitude.pdf、PDF物理11頁。

\textbf{公式正答：} 5（single）。

\par\textit{原典本文の校正と独自解説は未作成。分類のみの登録項目。}

\subsection{令和4年度 問3}
\label{q:r04-reg-apt-03}

\noindent\textbf{原典：} sources/R04/aptitude.pdf、PDF物理2頁。

\textbf{公式正答：} 3（single）。

\par\textit{原典本文の校正と独自解説は未作成。分類のみの登録項目。}
